\section{Introduction}
A common paradigm in theoretical computer science is to pin down which algorithmic techniques can or cannot achieve certain central goals by formalizing such a technique within a precise algorithmic model and then proving sharp limitations for that model. Classic examples include unconditional lower bounds for Linear Programming-based algorithms (i.e. extension complexity~\cite{FioriniMPTW15,Rothvoss17}), conditional lower bounds for preprocessing algorithms in parameterized complexity (conditioned on the non-collapse of the polynomial hierarchy~\cite{BodlaenderDFH09,FortnowS11}), and unconditional lower bounds for resolution-based algorithms for refuting the Strong Exponential Time Hypothesis~\cite{BeckI13}. This paradigm reveals which algorithmic nuances are fundamentally unavoidable and shows us where genuine breakthroughs must come from.

One notable example of such algorithmic nuance is that of algebraic cancellation: for several combinatorial computational problems, the fastest known algorithms crucially exploit certain algebraic cancellation in rather counter-intuitive ways, and it is a well-studied open question whether these algorithms can be matched with (more natural, and in some sense robust) `combinatorial' algorithms. This question arises for instance in the settings of Boolean matrix multiplication~\cite{WilliamsW18}, various variants of the problem of finding a perfect matching solvable by reduction to determinant computation~\cite{KarpUW86,MulmuleyVV87}, and fast exponential time algorithms for Hamiltonicity~\cite{Bjorklund14}. This very question of whether the use of algebraic cancellation is needed does not only occur in algorithm design, but also in various other disciplines. For example, in extremal combinatorics, it is an open question in various settings whether combinatorial proofs exist that match algebraic proofs (for example, a combinatorial proof of the skew two families theorem remains illusive~\cite{calbet2023k_r}), and the log rank conjecture in communication complexity asks whether algebraic low rank decompositions of matrices in general imply low rank combinatorial decompositions (or more specifically, whether the rank and partioning number of a Boolean matrix are quasi-polynomially related), see e.g.~\cite{LovettA14}.

One such setting in which we would like to replace algebraic arguments with combinatorial arguments within the area of parameterized complexity is that of connectivity problems parameterized by treewidth: In 2011 it was shown that that many connectivity problems parameterized by treewidth can be solved faster than what was deemed possible at the time~\cite{CyganNPPRW22}. In particular, problems like Hamiltonian Cycle and Steiner tree were solved with a randomized algorithm in $2^{O(k)}N^{O(1)}$ time, when given a tree decomposition of width $k$ of an $N$-vertex input graph. Before this work, fastest known algorithm for problems parameterized by treewidth were based on straightforward combinatorial dynamic programming techniques that naturally lead to running times like $k^{O(k)}N^{O(1)}$ and seemed hard to improve. Later work that followed up on~\cite{CyganNPPRW22} and provided deterministic algorithms for weighted extensions such as the Traveling Salesperson Problem (TSP) on bounded treewidth graphs~\cite{BodlaenderCKN15} and an alternative version connected the approach with a matroid-based extension of the aforementioned two-family theorem~\cite{FominLPS16}.

All these works~\cite{CyganNPPRW22,BodlaenderCKN15,FominLPS16} rely on algebraic cancellation by exploiting that the rank of certain \emph{compatibility} matrices is small, a particularly counterintuitive approach in comparison to the previous simple and clean combinatorial $k^{O(k)}N^{O(1)}$ time algorithms. It is a notable open question to remove the mentioned overhead in the runtime for the deterministic algorithm and weighted extensions. For example, a positive resolution of this open question would bring us closer to solving the TSP in time $1.9999^N$~\cite{Nederlof20}, where $N$ is the number of cities.

But, to do so it seems crucial to get a better combinatorial understanding of the involved compatibility matrices, simultaneously bypassing the undesired algebraic cancellation arguments. In this work we give a strong indication that more than just elementary combinatorial techniques are needed, via the algorithmic model of \emph{tropical circuits}.

\paragraph{Tropical Circuits.} A tropical circuit is an arithmetic circuit in which the inputs are labeled with variables that take an integer as value and the two arithmetic gates correspond to the max and sum, or alternatively, min and sum, operations. See \Cref{sec:prel} for a precise definition. The motivation for studying the expressiveness of tropical circuits is that it models a broad class dynamic programming algorithm. In particular, if there exists for a maximization problem (for which its instances are defined by a set of input weights) a dynamic programming algorithm that is `pure' in the sense that the associated recurrence only features the operations max and sum, than there is tropical circuit that outputs the optimal objective value with circuit size being proportional to the running time of the dynamic programming algorithm.

For example the classic Bellman-Held-Karp~\cite{Bellman62,heldKarp} $O(N^22^N)$ time algorithm for TSP is naturally converted into a tropical circuit with $O(N^22^N)$ gates and the $O(N^3)$ time Floyd-Warshall~\cite{Floyd62a,Warshall62} algorithm that computes the shortest path lengths between each pair of vertices is naturally converted into a tropical circuit with $O(N^3)$ gates. Both tropical circuits cannot be substantially improved~\cite{JerrumS82,kerr1970effect} (see also ~\cite[Corollary~2.2]{Jukna2023}).

In the realm of parameterized complexity, it is easy to see that canonical applications of dynamic programming can be modeled in an efficient way as tropical circuits. Examples are the dynamic programming algorithms for Steiner Tree and Set Cover (see e.g.~\cite[Section 6.1]{CyganFKLMPPS15}) with few number of terminals and elements and, especially relevant for this paper, the $O(2^{k} N)$ time algorithm for the maximum independent set problem on graphs with a given tree decomposition of width $k$ and $N$ vertices.

For much more detail on tropical circuits, we refer to the excellent textbook by Jukna~\cite{Jukna2023}.

\paragraph{Problems parameterized by treewidth.}
A very popular research line that started in~\cite{LokshtanovMS18b,LokshtanovMS18a} is that of investigating the fine-grained complexity of various NP-hard problems parameterized by width measures such as the treewidth of the input graph. 
In particular, for many NP-hard problems we are now able to design algorithms with a running time of the type $f(k)N^{O(1)}$ for some width measure $k$ and $n$ denoting the number of vertices of the input graph and simultaneously can prove that any improvement of this running time to $f(k)^{1-\Omega(1)}N^{O(1)}$ or even $f(k)^{o(1)}N^{O(1)}$ violates the Exponential Time Hypothesis (ETH) or the Strong Exponential Time Hypothesis (SETH). Such algorithms are often called \emph{(S)ETH-tight algorithms}. Such tight algorithms provide insight on how amenable the problem at hand is for divide and conquer algorithms since, conditioned on standard hypotheses, they reveal how much information of partial subsolutions exactly is relevant.

The aforementioned type of \emph{connectivity problems} such as TSP and Steiner Tree forms an important class of problems for which we do not generally have SETH-tight algorithms, precisely because the employed algebraic algorithms form a bottleneck towards deterministic algorithms and extensions to weighted variants that seemingly should be replaced with combinatorial arguments.


\paragraph{Our results.}
In this paper we provide evidence that direct combinatorial techniques on their own are insufficient for designing faster algorithm for connectivity problems, by giving unconditional lower bounds for tropical circuits. We state our lower bounds in terms of the~\emph{pathwidth} of the input graph. This is similar to the treewidth of a graph, except that we require more specifically to decompose the graph in a path-like manner instead of a tree-like manner. Hence, the pathwidth of a graph is always at least the treewidth of a graph and hence our lower bounds also imply lower bounds parameterized by treewidth. See \Cref{sec:prel} for definitions.
Before we study connectivity problems, we first study the complexity of a more basic problem:

\subparagraph{Maximum Weight Independent Set.}
In the Maximum Weight Independent Set problem one is given a graph $G=(V,E)$ along with a vertex weight $x_v$ for every $v \in V$, and is asked for the value $IS_G:=\max_{I \in \mathcal{I}(G)}\sum_{v \in I}x_v$ where $\mathcal{I}(G)$ denotes the family of all independent sets of $G$. Our lower bound for this problem reads as follows (full definitions are postponed to \Cref{sec:prel}):

\begin{theorem}\label{thm:main_is}
    For any $k \geq 1$, there exists a graph $G$ of pathwidth at most $k$ on $k^{O(1)}$ vertices such that any tropical circuit calculating $IS_G$ uses at least $\Omega(2^k)$ gates.
\end{theorem}
Note that pathwidth is a width measure that is always at least the treewidth of the graph.
Since there is a simple $O(2^{k}N)$-sized tropical circuit that calculates $IS_G$ for an $N$-vertex graph of pathwidth $k$, this result is optimal in a tight sense.
A previous result by Korhonen~\cite{Korhonen21} showed that for \emph{every} graph $G$ of treewidth $k$ and maximum degree $d$ any tropical circuit calculating $IS_G$ must be of size at least $2^{\Omega(k/d)}$.
This result is less tight than our new result and also seems to crucially rely on some properties of the Maximum Weight Independent Set problem. While the type of universal lower bound from~\cite{Korhonen21} is quite interesting, it does not directly have added value in our context of worst case complexity analysis and hence we do not pursue it further in this work.

\Cref{thm:main_is} is obtained using the following two ingredients. The graphs for which we show the bounds is constructed based on ideas from the classic reduction from CNF-Sat to Maximum Weight Independent Set. The tropical circuit size bound itself is shown by analyzing the structure of the so-called rectangles, a combinatorial notion that typically arises in the studie of tropical circuits that captures the way in which the partial solution calculated at some node of the tropical circuit can combine with the computations done by the rest of the circuit, see e.g.~\cite{Jukna2023}.

\subparagraph{Connectivity Problems.}
We show lower bounds on the tropical circuit complexity of the following graph connectivity problems. Let $G$ be an $N$-vertex graph with for every pair of distinct vertices $u$ and $v$ an edge weight $x_{u,v}\in \mathbb{N}$. Let $\mathcal{H}(G)$ denote the family of all sequences of $N + 1$ vertices $(u_i)_{i =0}^N$ such that $u_0 = u_N$ and $u_0 \to u_1 \to \dots \to u_{N - 1} \to u_N$ is a directed Hamiltonian cycle in $G$.
We define $DTSP_G$ to be the minimum $\sum_{i\in [N]}x_{u_{i-1},u_i}$ taken over all $(u_0,\ldots,u_N) \in \mathcal{H}(G)$. For an undirected graph, we define the undirected variant $TSP_G$ similarly while identifying variables $x_{u,v}$ and $x_{v, u}$ for all $uv \in E(G)$.

With these definitions in place, our main results can be stated as follows:
\begin{theorem}\label{thm:main_dtsp}
    For any $k \geq 1$, there exists a graph $G$ of pathwidth at most $k$ on $k^{O(1)}$ vertices such that any tropical circuit calculating $DTSP_G$ uses at least $2^{\Omega(k \log \log k)}$ gates.
\end{theorem}

\begin{theorem}\label{thm:main_tsp}
    For any $k \geq 1$, there exists a graph $G$ of pathwidth at most $k$ on $k^{O(1)}$ vertices such that any tropical circuit calculating $TSP_{G}$ uses at least $2^{\Omega(k \log \log k)}$ gates.
\end{theorem}
We also provide a similar lower bound for the Directed Spanning Tree problem.
We let $\mathcal{T}(G)$ denote the family of all functions $p : V(G) \to V(G)$ such that for exactly one vertex $v \in V(G)$, $p(v) = v$, and edges $(p(u), u)$ over all other vertices $u$ form an out-tree rooted at $v$ (i.e. a tree oriented away from the root). Furthermore, we define $DST_G$ as the minimum $\sum_{v \in V(G) : p(v) \neq v} x_{p(v), v}$ taken over all $p \in \mathcal{T}(G)$ (i.e., $DST_G$ is the minimum weight spanning tree of $G$).

\begin{theorem}\label{thm:main_dst}
    For any $k \geq 1$, there exists a graph $G$ of pathwidth at most $k$ on $k^{O(1)}$ vertices such that any tropical circuit calculating $DST_{G}$ uses at least $2^{\Omega(k \log \log k)}$ gates.
\end{theorem}

By adapting textbook dynamic programming algorithms, one can directly obtain tropical circuits for $DTSP_G,TSP_G$ and $DST_G$ of size $k^{O(k)} N$ if $G$ has $N$ vertices and pathwith/treewidth $k$. Thus there is still a gap between the lower bound and upper bound, and in fact this seems closely related to a similar gap in the area of communication complexity (see Section~\ref{sec:con} for more details on this). Nevertheless, our results show that the currently fastest $2^{O(k)}N$ time algorithms for TSP and Directed Steiner Tree (which generalizes DST) cannot be matched with merely pure dynamic programming.

The proofs of these three results combine the basic ingredients of the proof of \Cref{thm:main_is} with the properties of the \emph{Matchings Compatibility Matrix}, defined as follows. The rows and columns of this binary matrix are both indexed by perfect matchings of a bipartite graph, and an entry in the matrix indicates whether the union of these two perfect matchings form a Hamiltonian cycle. This matrix has already been studied in~\cite{raz1995log} in the context of the log-rank conjecture in communication complexity, where a lower bound on the non-deterministic communication complexity and an upper bound of its rank was given. The rank of (a slight variant of) this matrix also turned out to be important for studying the complexity of the Hamiltonian Cycle and TSP problem, both parameterized by the path/treewidth of the input graph and parameterized by the number of vertices~\cite{CurticapeanLN18,CyganKN18,Nederlof20}.

Our main technical contribution is that the lower bound on the non-deterministic communication complexity of this matrix from~\cite{raz1995log} can be used to obtain the above lower bounds on the size of tropical circuits.

\paragraph{Organization.}
This paper is organized as follows:
In \cref{sec:prel} we provide the necessary definitions and preliminary tools used in the remainder of the paper. \Cref{sec:mwis} presents the proof of \cref{thm:main_is}.
\cref{sec:match} presents the aforementioned matchings compatibility matrix and the required results about its structure.
In \cref{sec:tsp} we prove \cref{thm:main_dtsp} and \cref{thm:main_tsp}.
In \cref{sec:st} we prove \cref{thm:main_dst}, and we provide some concluding remarks in \cref{sec:con}.

